% --- Archon physics formalization source begin ---
% archon:physics
% archon:covers PhyXMiniProblems/problem_phyx_mini_0390.lean
% archon:source-report reports/phyx_mini/problem_phyx_mini_0390.source.json

\chapter{Physics problem phyx\_mini\_0390}
\label{ch:PhyXMiniProblems_problem_phyx_mini_0390}

\paragraph{Problem source.}
\#\# Physical scenario

Two piston/cylinder arrangements, A and B, have their gas chambers connected by a pipe, as shown in figure. The cross-sectional areas are \$A\_A = 75\$ cm\$\textasciicircum{}2\$ and \$A\_B = 25\$ cm\$\textasciicircum{}2\$, with the piston mass in A being \$m\_A = 25\$ kg. Assume an outside pressure of 100 kPa and standard gravitation.

\#\# Question

Find the mass \$m\_B\$ so that none of the pistons have to rest on the bottom.

\#\# Answer choices

- A: A: 490kg

- B: B: 8.33kg

- C: C: 154kg

- D: D: 10.2kg

\#\# Figure caption (auxiliary; use the image as primary evidence)

The image depicts two connected cylindrical containers, A and B, filled with a fluid. The containers are open at the top, with atmospheric pressure \textbackslash{}(P\_0\textbackslash{}) labeled above both. 

Container A is smaller in size and on the left side, while container B is larger and on the right. They are connected at the bottom by a horizontal U-shaped pipe filled with the same fluid.

The surface of the fluid in both containers is at the same horizontal level, suggesting hydrostatic equilibrium. An arrow labeled \textbackslash{}(g\textbackslash{}) points downward, indicating the direction of gravitational force. The pressure \textbackslash{}(P\_0\textbackslash{}) is shown acting on the surface of the fluid in both containers. 

This setup is typically used to illustrate principles of fluid statics, such as Pascal's Law or the concept of communicating vessels.

\#\# Dataset metadata

Laws of Thermodynamics; Physical Model Grounding Reasoning, Multi-Formula Reasoning

\paragraph{Recorded answer/context.}
B: B: 8.33kg

\paragraph{Figure/image path.}
/root/proposal\_for\_physic/science-mango/phyx\_mini\_run/phyx\_data/test\_image/390.png

\paragraph{Formalization target.}
create a compiling Lean file with sorry bodies at `PhyXMiniProblems/problem\_phyx\_mini\_0390.lean`. The Lean declarations must preserve the physical quantities, dimensions or dimensional roles, figure labels, governing-law hypotheses, and final relation expressed by this problem.
Use Mathlib/Physlib names found through LeanExplore where available. If a physics API is missing, introduce faithful local abstractions rather than scalar placeholder aliases.

\begin{theorem}[Physics formalization target]
\label{thm:physics:phyx_mini_0390:target}
\lean{PhyXMiniProblems.ProblemPhyXMini0390.problem_phyx_mini_0390}
\uses{def:physics:phyx-mini-0390:phyxminiproblems-problemphyxmini0390-massinkilograms, def:physics:phyx-mini-0390:phyxminiproblems-problemphyxmini0390-connectedpistonsetup, def:physics:phyx-mini-0390:phyxminiproblems-problemphyxmini0390-matchesconnectedpistonscenario, def:physics:phyx-mini-0390:phyxminiproblems-problemphyxmini0390-matchesproblemreadouts, def:physics:phyx-mini-0390:phyxminiproblems-problemphyxmini0390-usesstandardgravity, def:physics:phyx-mini-0390:phyxminiproblems-problemphyxmini0390-matchessuppliedconnectedpistonfigure, def:physics:phyx-mini-0390:phyxminiproblems-problemphyxmini0390-hasphysicalconnectedpistonparameters, def:physics:phyx-mini-0390:phyxminiproblems-problemphyxmini0390-neitherpistonrestsonbottom, def:physics:phyx-mini-0390:phyxminiproblems-problemphyxmini0390-satisfiesconnectedpistonlaws, def:physics:phyx-mini-0390:phyxminiproblems-problemphyxmini0390-displayedmasskilograms, def:physics:phyx-mini-0390:phyxminiproblems-problemphyxmini0390-roundstonearesthundredthkilogram, def:physics:phyx-mini-0390:phyxminiproblems-problemphyxmini0390-isuniqueclosestchoice}
The assigned autoformalize agent should translate this physics problem into Lean declarations in the covered file, with theorem and lemma proof bodies written as `by sorry`.
\end{theorem}
\begin{proof}
This is an autoformalization task, not a proof task. Produce statements that can later be proved by the physics prover without weakening the physical model.
\end{proof}

% --- Archon declaration topology begin ---
\section{Declaration topology}
\begin{definition}[Mass Quantity]
  \label{def:physics:phyx-mini-0390:phyxminiproblems-problemphyxmini0390-massquantity}
  \lean{PhyXMiniProblems.ProblemPhyXMini0390.MassQuantity}
  A nonnegative physical mass, independent of a choice of units
\end{definition}
\begin{proof}
  This is the declared model definition of Mass Quantity.
\end{proof}
\begin{definition}[Acceleration Quantity]
  \label{def:physics:phyx-mini-0390:phyxminiproblems-problemphyxmini0390-accelerationquantity}
  \lean{PhyXMiniProblems.ProblemPhyXMini0390.AccelerationQuantity}
  A nonnegative acceleration magnitude with physical dimension L T⁻²
\end{definition}
\begin{proof}
  This is the declared model definition of Acceleration Quantity.
\end{proof}
\begin{definition}[area In Square Meters]
  \label{def:physics:phyx-mini-0390:phyxminiproblems-problemphyxmini0390-areainsquaremeters}
  \lean{PhyXMiniProblems.ProblemPhyXMini0390.areaInSquareMeters}
  The numerical value of an area in square metres
\end{definition}
\begin{proof}
  This is the declared model definition of area In Square Meters.
\end{proof}
\begin{definition}[area In Square Centimeters]
  \label{def:physics:phyx-mini-0390:phyxminiproblems-problemphyxmini0390-areainsquarecentimeters}
  \lean{PhyXMiniProblems.ProblemPhyXMini0390.areaInSquareCentimeters}
  \uses{def:physics:phyx-mini-0390:phyxminiproblems-problemphyxmini0390-areainsquaremeters}
  The numerical value of an area in square centimetres
\end{definition}
\begin{proof}
  This is the declared model definition of area In Square Centimeters.
\end{proof}
\begin{definition}[mass In Kilograms]
  \label{def:physics:phyx-mini-0390:phyxminiproblems-problemphyxmini0390-massinkilograms}
  \lean{PhyXMiniProblems.ProblemPhyXMini0390.massInKilograms}
  \uses{def:physics:phyx-mini-0390:phyxminiproblems-problemphyxmini0390-massquantity}
  The numerical value of a mass in kilograms
\end{definition}
\begin{proof}
  This is the declared model definition of mass In Kilograms.
\end{proof}
\begin{definition}[acceleration In Meters Per Second Squared]
  \label{def:physics:phyx-mini-0390:phyxminiproblems-problemphyxmini0390-accelerationinmeterspersecondsquared}
  \lean{PhyXMiniProblems.ProblemPhyXMini0390.accelerationInMetersPerSecondSquared}
  \uses{def:physics:phyx-mini-0390:phyxminiproblems-problemphyxmini0390-accelerationquantity}
  The numerical value of an acceleration in metres per second squared
\end{definition}
\begin{proof}
  This is the declared model definition of acceleration In Meters Per Second Squared.
\end{proof}
\begin{definition}[pressure In Pascals]
  \label{def:physics:phyx-mini-0390:phyxminiproblems-problemphyxmini0390-pressureinpascals}
  \lean{PhyXMiniProblems.ProblemPhyXMini0390.pressureInPascals}
  The numerical value of an absolute pressure in pascals
\end{definition}
\begin{proof}
  This is the declared model definition of pressure In Pascals.
\end{proof}
\begin{definition}[pressure In Kilopascals]
  \label{def:physics:phyx-mini-0390:phyxminiproblems-problemphyxmini0390-pressureinkilopascals}
  \lean{PhyXMiniProblems.ProblemPhyXMini0390.pressureInKilopascals}
  \uses{def:physics:phyx-mini-0390:phyxminiproblems-problemphyxmini0390-pressureinpascals}
  The numerical value of an absolute pressure in kilopascals
\end{definition}
\begin{proof}
  This is the declared model definition of pressure In Kilopascals.
\end{proof}
\begin{definition}[Cylinder]
  \label{def:physics:phyx-mini-0390:phyxminiproblems-problemphyxmini0390-cylinder}
  \lean{PhyXMiniProblems.ProblemPhyXMini0390.Cylinder}
  The two piston/cylinder branches labelled in the figure
\end{definition}
\begin{proof}
  This is the declared model definition of Cylinder.
\end{proof}
\begin{definition}[Piston Support Regime]
  \label{def:physics:phyx-mini-0390:phyxminiproblems-problemphyxmini0390-pistonsupportregime}
  \lean{PhyXMiniProblems.ProblemPhyXMini0390.PistonSupportRegime}
  Whether a piston is freely supported or is supported by the cylinder base
\end{definition}
\begin{proof}
  This is the declared model definition of Piston Support Regime.
\end{proof}
\begin{definition}[Chamber Connection Model]
  \label{def:physics:phyx-mini-0390:phyxminiproblems-problemphyxmini0390-chamberconnectionmodel}
  \lean{PhyXMiniProblems.ProblemPhyXMini0390.ChamberConnectionModel}
  Whether the regions below the two pistons form one connected chamber
\end{definition}
\begin{proof}
  This is the declared model definition of Chamber Connection Model.
\end{proof}
\begin{definition}[Figure Label]
  \label{def:physics:phyx-mini-0390:phyxminiproblems-problemphyxmini0390-figurelabel}
  \lean{PhyXMiniProblems.ProblemPhyXMini0390.FigureLabel}
  Literal symbolic labels visible in the supplied figure
\end{definition}
\begin{proof}
  This is the declared model definition of Figure Label.
\end{proof}
\begin{definition}[Connected Piston Figure]
  \label{def:physics:phyx-mini-0390:phyxminiproblems-problemphyxmini0390-connectedpistonfigure}
  \lean{PhyXMiniProblems.ProblemPhyXMini0390.ConnectedPistonFigure}
  \uses{def:physics:phyx-mini-0390:phyxminiproblems-problemphyxmini0390-cylinder, def:physics:phyx-mini-0390:phyxminiproblems-problemphyxmini0390-figurelabel}
  Qualitative geometry and annotations read from the primary image
\end{definition}
\begin{proof}
  This is the declared model definition of Connected Piston Figure.
\end{proof}
\begin{definition}[Connected Piston Setup]
  \label{def:physics:phyx-mini-0390:phyxminiproblems-problemphyxmini0390-connectedpistonsetup}
  \lean{PhyXMiniProblems.ProblemPhyXMini0390.ConnectedPistonSetup}
  \uses{def:physics:phyx-mini-0390:phyxminiproblems-problemphyxmini0390-massquantity, def:physics:phyx-mini-0390:phyxminiproblems-problemphyxmini0390-accelerationquantity, def:physics:phyx-mini-0390:phyxminiproblems-problemphyxmini0390-cylinder, def:physics:phyx-mini-0390:phyxminiproblems-problemphyxmini0390-pistonsupportregime, def:physics:phyx-mini-0390:phyxminiproblems-problemphyxmini0390-chamberconnectionmodel, def:physics:phyx-mini-0390:phyxminiproblems-problemphyxmini0390-connectedpistonfigure}
  The physical parameters and observables of the apparatus. Pressures below the two pistons are kept separate so that equality of connected-chamber pressure is a governing law rather than a definitional shortcut. In particular, pistonMass .B is an independent physical quantity
\end{definition}
\begin{proof}
  This is the declared model definition of Connected Piston Setup.
\end{proof}
\begin{definition}[Matches Connected Piston Scenario]
  \label{def:physics:phyx-mini-0390:phyxminiproblems-problemphyxmini0390-matchesconnectedpistonscenario}
  \lean{PhyXMiniProblems.ProblemPhyXMini0390.MatchesConnectedPistonScenario}
  \uses{def:physics:phyx-mini-0390:phyxminiproblems-problemphyxmini0390-connectedpistonsetup}
  The chambers below A and B are connected by the pipe described in the prose
\end{definition}
\begin{proof}
  This is the declared model definition of Matches Connected Piston Scenario.
\end{proof}
\begin{definition}[Matches Problem Readouts]
  \label{def:physics:phyx-mini-0390:phyxminiproblems-problemphyxmini0390-matchesproblemreadouts}
  \lean{PhyXMiniProblems.ProblemPhyXMini0390.MatchesProblemReadouts}
  \uses{def:physics:phyx-mini-0390:phyxminiproblems-problemphyxmini0390-areainsquarecentimeters, def:physics:phyx-mini-0390:phyxminiproblems-problemphyxmini0390-massinkilograms, def:physics:phyx-mini-0390:phyxminiproblems-problemphyxmini0390-pressureinkilopascals, def:physics:phyx-mini-0390:phyxminiproblems-problemphyxmini0390-connectedpistonsetup}
  Numerical data supplied by the problem. The unknown mass of piston B is not a field of this structure
\end{definition}
\begin{proof}
  This is the declared model definition of Matches Problem Readouts.
\end{proof}
\begin{definition}[Uses Standard Gravity]
  \label{def:physics:phyx-mini-0390:phyxminiproblems-problemphyxmini0390-usesstandardgravity}
  \lean{PhyXMiniProblems.ProblemPhyXMini0390.UsesStandardGravity}
  \uses{def:physics:phyx-mini-0390:phyxminiproblems-problemphyxmini0390-accelerationinmeterspersecondsquared, def:physics:phyx-mini-0390:phyxminiproblems-problemphyxmini0390-connectedpistonsetup}
  The exact conventional calibration g₀ = 9.80665 m/s²
\end{definition}
\begin{proof}
  This is the declared model definition of Uses Standard Gravity.
\end{proof}
\begin{definition}[Matches Supplied Connected Piston Figure]
  \label{def:physics:phyx-mini-0390:phyxminiproblems-problemphyxmini0390-matchessuppliedconnectedpistonfigure}
  \lean{PhyXMiniProblems.ProblemPhyXMini0390.MatchesSuppliedConnectedPistonFigure}
  \uses{def:physics:phyx-mini-0390:phyxminiproblems-problemphyxmini0390-connectedpistonsetup}
  Primary-image evidence. The image supplies geometry and labels but no value for the mass of piston B
\end{definition}
\begin{proof}
  This is the declared model definition of Matches Supplied Connected Piston Figure.
\end{proof}
\begin{definition}[Has Physical Connected Piston Parameters]
  \label{def:physics:phyx-mini-0390:phyxminiproblems-problemphyxmini0390-hasphysicalconnectedpistonparameters}
  \lean{PhyXMiniProblems.ProblemPhyXMini0390.HasPhysicalConnectedPistonParameters}
  \uses{def:physics:phyx-mini-0390:phyxminiproblems-problemphyxmini0390-areainsquaremeters, def:physics:phyx-mini-0390:phyxminiproblems-problemphyxmini0390-massinkilograms, def:physics:phyx-mini-0390:phyxminiproblems-problemphyxmini0390-accelerationinmeterspersecondsquared, def:physics:phyx-mini-0390:phyxminiproblems-problemphyxmini0390-pressureinpascals, def:physics:phyx-mini-0390:phyxminiproblems-problemphyxmini0390-connectedpistonsetup}
  Positivity and nondegeneracy conditions for the physical apparatus
\end{definition}
\begin{proof}
  This is the declared model definition of Has Physical Connected Piston Parameters.
\end{proof}
\begin{definition}[Neither Piston Rests On Bottom]
  \label{def:physics:phyx-mini-0390:phyxminiproblems-problemphyxmini0390-neitherpistonrestsonbottom}
  \lean{PhyXMiniProblems.ProblemPhyXMini0390.NeitherPistonRestsOnBottom}
  \uses{def:physics:phyx-mini-0390:phyxminiproblems-problemphyxmini0390-connectedpistonsetup}
  The operating condition in the question: neither piston receives a support reaction from the bottom of its cylinder
\end{definition}
\begin{proof}
  This is the declared model definition of Neither Piston Rests On Bottom.
\end{proof}
\begin{definition}[Satisfies Connected Piston Laws]
  \label{def:physics:phyx-mini-0390:phyxminiproblems-problemphyxmini0390-satisfiesconnectedpistonlaws}
  \lean{PhyXMiniProblems.ProblemPhyXMini0390.SatisfiesConnectedPistonLaws}
  \uses{def:physics:phyx-mini-0390:phyxminiproblems-problemphyxmini0390-areainsquaremeters, def:physics:phyx-mini-0390:phyxminiproblems-problemphyxmini0390-massinkilograms, def:physics:phyx-mini-0390:phyxminiproblems-problemphyxmini0390-accelerationinmeterspersecondsquared, def:physics:phyx-mini-0390:phyxminiproblems-problemphyxmini0390-pressureinpascals, def:physics:phyx-mini-0390:phyxminiproblems-problemphyxmini0390-connectedpistonsetup}
  The connected-pressure law and vertical force balance for each freely supported piston. In coherent SI readouts, force balance is P\_gas A = P₀ A + m g. It is conditional on free support because a piston resting on the base would also have a bottom reaction force. Neither field states the desired mass ratio or the numerical mass of piston B
\end{definition}
\begin{proof}
  This is the declared model definition of Satisfies Connected Piston Laws.
\end{proof}
\begin{theorem}[freely Supported Piston Mass Ratio]
  \label{thm:physics:phyx-mini-0390:phyxminiproblems-problemphyxmini0390-freelysupportedpistonmassratio}
  \lean{PhyXMiniProblems.ProblemPhyXMini0390.freelySupportedPistonMassRatio}
  \uses{def:physics:phyx-mini-0390:phyxminiproblems-problemphyxmini0390-areainsquaremeters, def:physics:phyx-mini-0390:phyxminiproblems-problemphyxmini0390-massinkilograms, def:physics:phyx-mini-0390:phyxminiproblems-problemphyxmini0390-connectedpistonsetup, def:physics:phyx-mini-0390:phyxminiproblems-problemphyxmini0390-matchesconnectedpistonscenario, def:physics:phyx-mini-0390:phyxminiproblems-problemphyxmini0390-hasphysicalconnectedpistonparameters, def:physics:phyx-mini-0390:phyxminiproblems-problemphyxmini0390-neitherpistonrestsonbottom, def:physics:phyx-mini-0390:phyxminiproblems-problemphyxmini0390-satisfiesconnectedpistonlaws}
  Equal pressure and force balance make piston mass proportional to area
\end{theorem}
\begin{proof}
  The conclusion follows from the governing relations and figure readouts named in the statement of freely Supported Piston Mass Ratio.
\end{proof}
\begin{theorem}[piston BMass Is Twenty Five Thirds Kilograms]
  \label{thm:physics:phyx-mini-0390:phyxminiproblems-problemphyxmini0390-pistonbmassistwentyfivethirdskilograms}
  \lean{PhyXMiniProblems.ProblemPhyXMini0390.pistonBMassIsTwentyFiveThirdsKilograms}
  \uses{def:physics:phyx-mini-0390:phyxminiproblems-problemphyxmini0390-massinkilograms, def:physics:phyx-mini-0390:phyxminiproblems-problemphyxmini0390-connectedpistonsetup, def:physics:phyx-mini-0390:phyxminiproblems-problemphyxmini0390-matchesconnectedpistonscenario, def:physics:phyx-mini-0390:phyxminiproblems-problemphyxmini0390-matchesproblemreadouts, def:physics:phyx-mini-0390:phyxminiproblems-problemphyxmini0390-hasphysicalconnectedpistonparameters, def:physics:phyx-mini-0390:phyxminiproblems-problemphyxmini0390-neitherpistonrestsonbottom, def:physics:phyx-mini-0390:phyxminiproblems-problemphyxmini0390-satisfiesconnectedpistonlaws}
  Substitution of the two areas and mass A gives the exact mass of piston B. The outside pressure and gravitational acceleration cancel from the relation
\end{theorem}
\begin{proof}
  The conclusion follows from the governing relations and figure readouts named in the statement of piston BMass Is Twenty Five Thirds Kilograms.
\end{proof}
\begin{definition}[Answer Choice]
  \label{def:physics:phyx-mini-0390:phyxminiproblems-problemphyxmini0390-answerchoice}
  \lean{PhyXMiniProblems.ProblemPhyXMini0390.AnswerChoice}
  Labels of the four answer choices in the source problem
\end{definition}
\begin{proof}
  This is the declared model definition of Answer Choice.
\end{proof}
\begin{definition}[displayed Mass Kilograms]
  \label{def:physics:phyx-mini-0390:phyxminiproblems-problemphyxmini0390-displayedmasskilograms}
  \lean{PhyXMiniProblems.ProblemPhyXMini0390.displayedMassKilograms}
  \uses{def:physics:phyx-mini-0390:phyxminiproblems-problemphyxmini0390-answerchoice}
  Piston-B mass in kilograms displayed beside each answer label
\end{definition}
\begin{proof}
  This is the declared model definition of displayed Mass Kilograms.
\end{proof}
\begin{definition}[recorded Dataset Answer]
  \label{def:physics:phyx-mini-0390:phyxminiproblems-problemphyxmini0390-recordeddatasetanswer}
  \lean{PhyXMiniProblems.ProblemPhyXMini0390.recordedDatasetAnswer}
  \uses{def:physics:phyx-mini-0390:phyxminiproblems-problemphyxmini0390-answerchoice}
  Dataset metadata records answer label B
\end{definition}
\begin{proof}
  This is the declared model definition of recorded Dataset Answer.
\end{proof}
\begin{definition}[Rounds To Nearest Hundredth Kilogram]
  \label{def:physics:phyx-mini-0390:phyxminiproblems-problemphyxmini0390-roundstonearesthundredthkilogram}
  \lean{PhyXMiniProblems.ProblemPhyXMini0390.RoundsToNearestHundredthKilogram}
  An exact mass rounds to a displayed value to the nearest hundredth of a kilogram when the error is less than half of 0.01 kg
\end{definition}
\begin{proof}
  This is the declared model definition of Rounds To Nearest Hundredth Kilogram.
\end{proof}
\begin{definition}[Is Unique Closest Choice]
  \label{def:physics:phyx-mini-0390:phyxminiproblems-problemphyxmini0390-isuniqueclosestchoice}
  \lean{PhyXMiniProblems.ProblemPhyXMini0390.IsUniqueClosestChoice}
  \uses{def:physics:phyx-mini-0390:phyxminiproblems-problemphyxmini0390-answerchoice, def:physics:phyx-mini-0390:phyxminiproblems-problemphyxmini0390-displayedmasskilograms}
  A chosen answer is strictly closer to the exact mass than every rival
\end{definition}
\begin{proof}
  This is the declared model definition of Is Unique Closest Choice.
\end{proof}
% --- Archon declaration topology end ---
% --- Archon physics formalization source end ---
